\section{The closure ordinal and the strict-finite regime}

This section records claim nodes and dependency edges.

\begin{proposition}[The settling orbit $f^n$]\label{thm:iter}\lean{Obsmetrologia.Charge.iter}\leanok\uses{thm:closure-operator}
$\texttt{iter}\,f\,n\,x = f^n(x)$, plain recursion iterating an endomap $n$ times
from $x$. This is the orbit whose stabilization stage is the closure ordinal.
\end{proposition}
\begin{theorem}[Fixed points stay fixed]\label{thm:iter-fixed-of}\lean{Obsmetrologia.Charge.iter_fixed_of}\leanok\uses{thm:iter}
If $f\,a = a$ then $\texttt{iter}\,f\,n\,a = a$ for all $n$: once the orbit
reaches a fixed point it never leaves it.
\end{theorem}
\begin{theorem}[Iteration adds]\label{thm:iter-add}\lean{Obsmetrologia.Charge.iter_add}\leanok\uses{thm:iter}
$f^{n+m}(x) = f^m(f^n(x))$, the additive law of iteration, by induction on $m$.
\end{theorem}
\begin{theorem}[The orbit is a $\le_k$-chain]\label{thm:iter-step-le}\lean{Obsmetrologia.Charge.iter_step_le}\leanok\uses{thm:iter, thm:kbot-ktop}
For an inflationary $f$, each step climbs: $f^n(x) \le_k f^{n+1}(x)$. The
settling orbit is therefore an ascending knowledge chain.
\end{theorem}
\begin{theorem}[No ascending chain of four distinct charges]\label{thm:kchain-height}\lean{Obsmetrologia.Charge.kchain_height}\leanok\uses{thm:kle}
For any $a \le_k b \le_k c \le_k d$, at least one of $a=b$, $b=c$, $c=d$ holds:
the knowledge order has height two, so a chain of four must collapse somewhere.
A finite \texttt{decide} check; this is the pigeonhole that bounds the closure
ordinal.
\end{theorem}
\begin{theorem}[Propagation from coincidence]\label{thm:iter-propagate}\lean{Obsmetrologia.Charge.iter_propagate}\leanok\uses{thm:iter-add, thm:iter-fixed-of}
If two consecutive iterates coincide at stage $k$ then the orbit is constant
from $k$ on: $\texttt{iter}\,f\,(k+j)\,x = \texttt{iter}\,f\,k\,x$ for all $j$.
Once it stops, it stays stopped.
\end{theorem}
\begin{theorem}[Strict-finite termination: $f^3 = f^4$]\label{thm:settles-by-three}\lean{Obsmetrologia.Charge.settles_by_three}\leanok\uses{thm:iter-step-le, thm:kchain-height, thm:iter-propagate}
Every inflationary endomap on the charge lattice reaches its fixed point by step
three: $\texttt{iter}\,f\,3\,x = \texttt{iter}\,f\,4\,x$. The height-two chain
bound forces one of the first three steps to be stationary, and propagation
carries it forward. The closure ordinal is $\le 3$.
\end{theorem}
\begin{proposition}[The settling bound $\texttt{SettlesBy}$]\label{thm:settles-by}\lean{Obsmetrologia.Charge.SettlesBy}\leanok\uses{thm:iter}
$\texttt{SettlesBy}\,f\,x\,n$ holds when $f^{n+1}(x) = f^n(x)$ --- the orbit is
fixed from stage $n$. The least such $n$ is the closure ordinal $\kappa$; any
such $n$ is a settling bound $\lambda$.
\end{proposition}
\begin{theorem}[The strict-finite regime: $\kappa \le 3$]\label{thm:strict-finite}\lean{Obsmetrologia.Charge.strict_finite}\leanok\uses{thm:settles-by, thm:settles-by-three}
Every inflationary endomap on the charge lattice satisfies
$\texttt{SettlesBy}\,f\,x\,3$ --- a uniform finite settling bound with no
transfinite tail. This is the strict-finite regime named in the module
docstring.
\end{theorem}
\begin{proposition}[The four closure-ordinal regimes]\label{thm:closure-regime-carrier}\lean{Obsmetrologia.ClosureRegime}\leanok\uses{thm:strict-finite, thm:corner}
The closure ordinal does not produce an arbitrary taxonomy.  It lands in the
recurring FOUR as the regime carrier: strict, Scott, transfinite-bounded, and
aperiodic.
\end{proposition}
\begin{proposition}[Regime data are the decidable regime coordinates]\label{thm:closure-regime-data}\lean{Obsmetrologia.ClosureRegimeData}\leanok\uses{thm:closure-regime-carrier}
A settling problem is projected to three decidable coordinates: Scott
continuity, finite base, and countable cofinality, with the finite-base
constraints carried as fields.
\end{proposition}
\begin{theorem}[The four-regime dichotomy]\label{thm:four-regime-dichotomy}\lean{Obsmetrologia.four_regime_dichotomy}\leanok\uses{thm:closure-regime-data}
Every regime datum selects exactly the strict, Scott, transfinite-bounded, or
aperiodic case, with the intended Boolean characterization for each case.
\end{theorem}
\begin{theorem}[The regime carrier is the FOUR again]\label{thm:regime-is-four-corner}\lean{Obsmetrologia.regime_is_four_corner}\leanok\uses{thm:four-regime-dichotomy, thm:corner}
The regime carrier round-trips with the articulation corners.  This is a
specific incidence weld, not a claim that the regime chain is already the same
ordered object as the knowledge diamond.
\end{theorem}
\begin{theorem}[The spine classifies as strict]\label{thm:spine-is-strict-regime}\lean{Obsmetrologia.spine_is_strict}\leanok\uses{thm:strict-finite, thm:four-regime-dichotomy}
The finite charge-lattice spine has finite-base regime data, so its regime is
strict.  The strict-finite theorem is the operational witness.
\end{theorem}
\begin{theorem}[Regime cohomology vanishes exactly on strict/Scott]\label{thm:regime-cohomology-vanishes-iff}\lean{Obsmetrologia.cohomologyVanishes_iff}\leanok\uses{thm:four-regime-dichotomy, thm:spine-is-strict-regime}
The regime table marks the flat side: cohomology vanishes exactly for strict
and Scott regimes, and can carry content on the transfinite-bounded and
aperiodic side.
\end{theorem}
\begin{theorem}[The settled value is genuinely fixed]\label{thm:settled-fixed}\lean{Obsmetrologia.Charge.settled_fixed}\leanok\uses{thm:settles-by-three}
$f(f^3(x)) = f^3(x)$: after three steps the orbit is at an honest fixed point,
not merely slowing.
\end{theorem}
\begin{corollary}[Every closure settles in $\le 3$ steps]\label{cor:closure-strict-finite}\lean{Obsmetrologia.Charge.closure_strict_finite}\leanok\uses{thm:closure-operator, thm:strict-finite}
Since every \texttt{Closure} is inflationary, $\texttt{SettlesBy}\,c.\texttt{op}\,x\,3$
holds --- consistent with idempotent closures settling in one step. The two
models agree: the one-step idempotent settling is the degenerate case of the
$\kappa \le 3$ iterated settling.
\end{corollary}
